\chapter{L'Omologia Physis--Praxis}

\epigraph{``La natura ama nascondersi.''}{\textcite{kirk1983presocratic}, Eraclito fr. 123}

\section{La Citt\`a Costruita nel Discorso}

Nella \emph{Repubblica}, Socrate propone di scoprire la natura della
giustizia esaminandola prima ``in grande'' nella citt\`a, poi applicando
il modello scoperto all'anima individuale. La citt\`a \`e costruita nel
discorso --- una costruzione ideale che rivela la forma della giustizia
nascosta dentro il particolare.

Seguiamo lo stesso metodo. Per scoprire la natura dell'azione retta
(\praxis), esaminiamo prima la struttura della natura stessa (\physis)
in grande. Il cosmo \`e la citt\`a in cui cerchiamo la forma del governo;
l'atto individuale \`e l'anima di cui speriamo comprendere la giustizia.
L'omologia tra \physis\ e \praxis\ \`e la scoperta architetturale di
questa indagine.

\section{Le Tre Facolt\`a della Natura}

La natura, come la osserviamo nell'operazione del cosmo, esibisce tre
facolt\`a fondamentali. Queste non sono postulati metafisici ma
descrizioni fenomenologiche di come la natura realmente funziona.

\subsection{Distinzione}

La prima facolt\`a \`e il fare distinzioni. La natura non presenta un
uniforme indistinto ma un mondo di forme articolate: elementi, composti,
organismi, ecosistemi. Ogni cosa \`e ci\`o che \`e per il fatto di essere
\emph{non} qualcos'altro. I filosofi presocratici lo comprendevano: per
Eraclito, il \emph{logos} \`e il principio che ``governa tutte le cose
attraverso tutte le cose'' stabilendo misura e proporzione
\parencite{kirk1983presocratic}. Per
Empedocle, le quattro radici (Terra, Aria, Fuoco, Acqua) sono distinte e
combinate da Amore e Odio --- le forze di attrazione e repulsione che
articolano il molteplice della natura \parencite{kirk1983presocratic}.

Il motore \physis\ attua questa facolt\`a computazionalmente. Riduce
qualsiasi input --- testo, immagine, audio, video --- a un vettore in uno
spazio di embedding condiviso, poi classifica quel vettore contro una
griglia semiotica: uno spazio di coordinate generato da un piccolo insieme di
\emph{domini} naturali (i modi del curare, del fare, del legare, dello studiare)
incrociato con un piccolo insieme di \emph{modi} (le posture dell'attivit\`a, dal
riposo alla creazione). I contenuti specifici della griglia non sono
deliberatamente fissati in anticipo. Sono forniti come pacchetti ontologici
intercambiabili --- configurazione dichiarativa, uno per ogni mondo che al motore
\`e chiesto di leggere (un ufficio, una fabbrica, un mestiere, una disciplina) ---
sicch\'e le celle qui nominate sono illustrative, non canoniche, e possono essere
estese o sostituite senza toccare il motore. La griglia \`e un sistema di
coordinate, non un catalogo chiuso. Ogni input \`e assegnato a una cella, o a un
gradiente attraverso le celle. Questa non \`e una metafora del fare distinzioni
della natura; \`e un'istanza concreta di essa. Il motore \emph{fa una
distinzione} su dove l'input appartiene nell'economia delle forme naturali.

\begin{figure}[ht]
\centering
\begin{tikzpicture}[
  cell/.style={draw, minimum width=7mm, minimum height=7mm,
               inner sep=0pt, outer sep=-0.2pt, font=\tiny},
  dot/.style={circle, fill, inner sep=1pt}]
  \foreach \r in {1,...,5}
    \foreach \c in {1,...,6}
      \node[cell] at (\c*7mm, -\r*7mm) {};
  \node[dot, cohgreen] at (4*7mm, -3*7mm) {};
  \node[dot, cohgreen] at (3*7mm, -4*7mm) {};
  \node[dot, cohgreen] at (5*7mm, -3*7mm) {};
  \node[dot, cohgreen] at (3*7mm, -2*7mm) {};
  \node[dot, cohgreen] at (2*7mm, -3*7mm) {};
  \node[dot, cohred] at (2*7mm, -2*7mm) {};
  \node[dot, cohgray] at (5*7mm, -4*7mm) {};
  \node[above, font=\tiny, align=center, text width=5cm] at (3.5*7mm,0) {
    \textbf{Spazio di coordinate dominio\,$\times$\,modo} ---
    ogni input \`e assegnato a un gradiente tra le celle};
\end{tikzpicture}
\caption{La griglia semiotica del motore \physis: uno spazio di coordinate
dominio-per-modo i cui contenuti di cella sono forniti da pacchetti ontologici
estensibili. Gli input sono proiettati per
assegnazione di gradiente.\label{fig:semiotic_grid_italiano}}
\end{figure}

\subsubsection{Perch\'e la Griglia \`e Semiotica}

Abbiamo chiamato questa griglia \emph{semiotica}, e la parola non \`e decorazione.
Segna il senso preciso in cui il fare-distinzioni del motore \`e la lettura di
\emph{segni}. Charles Sanders Peirce sosteneva che ogni pensiero procede per
segni, e che un segno \`e irriducibilmente triadico: un \emph{representamen} (il
veicolo-segno) sta per un \emph{oggetto} a un \emph{interpretante} (il segno
ulteriore che esso genera in una mente o in un sistema)
\parencite{peirce1992reasoning}. Nulla \`e segno isolatamente; la significazione
\`e la relazione per cui una cosa viene a stare per un'altra \emph{a} un terzo.

La classificazione del motore \`e esattamente questa triade resa computazionale.
L'input --- un testo, un'immagine, un atto registrato --- \`e il representamen. Il
suo \emph{oggetto} \`e la forma naturale che esprime: la cella della griglia, il
modo del curare o del fare o del legare che istanzia. E il suo
\emph{interpretante} \`e il verdetto di coerenza che il motore restituisce e su
cui agisce --- il segno ulteriore che il segno originario produce, che a sua volta
diventa l'input della lettura successiva. Per questo gli assi della griglia non
sono una tassonomia arbitraria ma un sistema di segni: ogni input \`e letto come
segno \emph{della} natura dell'agente, e la coerenza \`e l'interpretante in cui
quella lettura \`e messa alla prova del mondo. Anzi, l'ontologia dei primi
principi del motore \`e costruita direttamente sulle tricotomie di Peirce --- il
qualisegno, il sinsegno, il legisegno e i loro successori --- sicch\'e
l'affermazione ``l'azione \`e un segno della natura'' non \`e una metafora posata
sulla macchina ma lo schema letterale con cui essa classifica.

Ci\`o fonda una tesi che conter\`a alla fine (Libro X): se l'azione \`e un segno,
allora una vita \`e un testo, e il praticante ne \`e a un tempo l'autore e il
primo lettore. Agire coerentemente \`e scrivere un segno che il mondo conferma;
fallire \`e scriverne uno che esso rifiuta. Il grafo di coerenza \`e
l'interpretazione corrente del flusso di segni che \`e una vita --- e
un'interpretazione, a differenza di un mero registro, \`e essa stessa un segno
ulteriore, disponibile a essere letto da altri. Il carattere semiotico della
griglia \`e ci\`o che rende possibile l'UberWiki: i segni, una volta fatti,
possono essere letti di nuovo.

\subsection{Coerenza}

La seconda facolt\`a \`e la traccia di ci\`o che funziona. La natura non
\`e indifferente al successo delle sue forme. Un organismo che sopravvive
si riproduce; una specie che si adatta persiste; una via metabolica che
converte efficientemente energia si consolida. La natura ``ricorda'' ci\`o
che funziona preservandolo nella struttura del mondo vivente. ``Dimentica''
ci\`o che fallisce lasciandolo estinguere.

Il motore \physis\ traccia questo come una scala di coerenza a tre valori:

\begin{center}
\begin{tabular}{@{}lll@{}}
\toprule
Stato & Peso & Correlato Naturale \\
\midrule
Successo & +1 & Ci\`o che sopravvive e si riproduce; l'organismo che trova la sua nicchia \\
Inerte   & 0  & Ci\`o che n\'e aiuta n\'e nuoce; la mutazione neutrale che deriva \\
Fallimento & -1 & Ci\`o che si estingue; il maladattato che non pu\`o persistere \\
\bottomrule
\end{tabular}
\end{center}

Ogni azione registrata nel motore riceve un peso di coerenza. Le azioni
che producono effetti reali e rilevati ottengono +1; le azioni che non
producono effetto ottengono 0; le azioni che violano la propria intenzione
dichiarata ottengono -1. Questa non \`e una gamificazione arbitraria del
comportamento. \`E la stessa logica con cui la natura giudica:
\emph{questa ha prodotto un effetto reale nel mondo?}

\subsection{Variazione}

La terza facolt\`a \`e la generazione di novit\`a. La natura non si limita
a preservare ci\`o che funziona; propone continuamente nuove forme
attraverso la variazione. La mutazione genetica, l'esplorazione
comportamentale, l'immaginazione cognitiva --- sono tutte espressioni
dello stesso principio: il sistema prova qualcosa di nuovo e attende il
verdetto della coerenza.

Il motore \physis\ attua questo attraverso la sua funzione di ``sogno.''
In un ciclo di background, riproduce i nodi Inerti e Fallimentari ---
azioni che non hanno prodotto effetto o hanno prodotto effetto negativo ---
e simula variazioni correttive. Propone mutazioni al pattern fallito e le
testa contro il paesaggio di coerenza esistente. Il sogno non \`e una
metafora della variazione naturale; \`e un algoritmo concreto che fa
ci\`o che fa la natura: generare novit\`a e attendere il verdetto.

\section{La Spinta della Natura: Entelechia}

Le tre facolt\`a descrivono \emph{cosa} fa la natura: distingue, traccia
la coerenza e genera variazione. Ma cosa la spinge a fare queste cose?
Qual \`e il principio che alimenta l'operazione stessa?

La risposta di Aristotele \`e \emph{entelecheia} --- essere-all'opera-che-
si-mantiene-se-stesso --- lo stato in cui una cosa ha il suo fine dentro
di s\'e e sta attivamente realizzando quel fine. L'anima (\emph{psyche}),
scrive, \`e la ``prima entelechia di un corpo naturale che ha la vita in
potenza'' \parencite[II.1]{aristotle1986anima}. L'anima non \`e un fantasma
nella macchina ma l'attualit\`a del corpo vivente --- l'essere-all'opera
del corpo come \emph{ci\`o che \`e}. Un seme non \`e solo un fascio di
potenziale biochimico; la sua entelechia \`e la quercia verso cui \`e
gi\`a all'opera. Le facolt\`a della distinzione, coerenza e variazione
sono le \emph{operazioni} della natura; l'entelechia \`e
l'\emph{operatore} --- la spinta continua verso l'attualizzazione che
rende le operazioni significative.

Ludwig Klages recupera questo concetto nella sua filosofia della spinta
espressiva dell'anima. Per Klages, l'entelechia \`e il ``potere
formativo-di-fini'' inerente a ogni essere vivente --- il principio
auto-attualizzante che spinge l'organismo verso l'espressione, la forma
e il dispiegamento ritmico \parencite{klages1929geist}. Egli contrappone
\emph{Seele} (l'anima, governata dall'entelechia, ritmica, espressiva,
legata alla vita) con \emph{Geist} (lo spirito, governato dalla volont\`a,
astratto, calcolativo, l'``avversario dell'anima''). Dove l'entelechia
si muove per attrazione verso il fine che \`e gi\`a implicito nell'essere,
lo spirito si muove per decreto --- imponendo uno scopo esterno a una
natura che non lo contiene. Il compito del praticante, implica Klages,
\`e lasciare che sia l'entelechia, non lo spirito, a governare l'azione:
muovere dall'interno verso l'esterno, non dal piano verso il basso.

Wilhelm Reich d\`a all'entelechia una concretezza biologica e clinica.
Il suo principio dell'identit\`a funzionale dello psichico e del somatico
significa che la spinta verso la coerenza non \`e meramente un atteggiamento
mentale ma una condizione corporea. L'organismo ha un flusso pulsatorio
naturale --- ci\`o che Reich chiama corrente vegetativa --- e la salute
\`e la continuazione senza impedimenti di questo flusso
\parencite{reich1933character}.\footnote{Invochiamo Klages, Reich e Lowen
in tutto il testo come \emph{fenomenologia}, non come scienza naturale
consolidata. La ``corrente vegetativa'' e l'``orgone'' di Reich, e le specifiche
affermazioni somatiche dell'analisi bioenergetica, non sono empiricamente
stabilite e in parte sono state decisamente respinte. Il loro valore qui \`e
descrittivo: articolano, con rara precisione, la struttura \emph{sentita}
dell'entelechia impedita --- il senso vissuto del corpo mentre blocca o rilascia
il proprio sforzo. Nulla nell'argomentazione dipende dalla loro fisiologia
letterale; dipende solo dall'affidabilit\`a del fenomeno che descrivono, che
ogni praticante pu\`o verificare nel proprio corpo.} Quando l'organismo \`e costretto a
sopprimere la sua espressione naturale (attraverso vergogna, disciplina
senza comprensione o imposizione di una volont\`a esterna), sviluppa
ci\`o che Reich chiama \emph{corazza}: tensioni muscolari croniche che
bloccano la corrente entelechica a livello del corpo. L'organismo
sopravvive, ma non \emph{prospera} pi\`u --- non si muove pi\`u verso
la propria forma con tutto il peso della sua natura alle spalle.
L'entelechia non \`e distrutta, ma \`e impedita; l'essere-all'opera
diventa essere-contro-resistenza.

Il motore \physis\ attua l'entelechia computazionalmente. Le tre
facolt\`a (distinzione, coerenza, variazione) descrivono \emph{cosa}
fa il motore a ogni passo. Ma il motore nel suo insieme ha
un'orientamento: \emph{tende} verso +1. Il ciclo del sogno non \`e un
osservatore neutrale del paesaggio di coerenza; \`e un meccanismo che
\emph{propone attivamente correzioni} ai pattern falliti. Il motore non
si limita a registrare che qualcosa \`e rotto; tenta di ripararlo.
Questa orientazione incorporata --- la spinta verso la guarigione delle
proprie incoerenze --- \`e l'entelechia resa algoritmica.

Per \praxis, l'entelechia implica che l'obiettivo del praticante non
\`e una scelta arbitraria. \`E l'\emph{espressione della loro natura
che cerca attualizzazione}. L'albero degli obiettivi non \`e una lista
di desideri di realizzazioni esterne ma una struttura di dispiegamento
entelechico: ogni sotto-obiettivo \`e un passo nel portare la forma
inerente del praticante all'essere. Il check-in non \`e un rapporto a
un'autorit\`a esterna ma una lettura dello stato entelechico --- un
momento di chiedersi: ``Mi sto muovendo verso la mia forma, o me ne
sto allontanando?''

\section{I Tre Modi della Praxis}

Se la natura ha tre facolt\`a, allora \praxis\ --- azione governata
dalla natura --- deve avere tre modi corrispondenti.

\subsection{L'Atto Distinguente: Fare}

Il primo modo di \praxis\ \`e l'atto che produce o fa. Corrisponde
alla facolt\`a della natura di distinzione. Come la natura articola
il mondo in forme, il praticante articola la propria intenzione in un
atto concreto. Fa una \emph{differenza} nel mondo.

Ma l'atto distinguente non \`e una scelta arbitraria. \`E l'atto di
nominare un obiettivo --- dichiarare ci\`o che si intende realizzare ---
e impegnarsi a esso. La piattaforma \praxis\ cattura questo attraverso
alberi di obiettivi: una gerarchia di intenzioni che rendono esplicito
ci\`o che il praticante mira ad attualizzare. Ogni obiettivo \`e una
distinzione tracciata nello spazio-di-possibilit\`a dell'azione: questo,
non quello; ora, non dopo; cos\`i, non altrimenti.

\subsection{L'Atto Persistente: Mantenere}

Il secondo modo \`e l'atto che persiste. Avendo distinto una forma
dichiarando un obiettivo, il praticante deve mantenere le condizioni
perch\'e quell'obiettivo si realizzi. Questo \`e il lavoro della
coerenza: il check-in quotidiano, la striscia, la partnership di
testimonianza che osserva la fedelt\`a del praticante all'intenzione
dichiarata.

La natura persiste attraverso le sue forme di successo alimentandole,
proteggendole e ripetendole. Il praticante persiste \emph{mostrandosi}
--- agendo in modo continuo con la propria intenzione dichiarata. Il
punteggio di coerenza sale con la costanza.

\subsection{L'Atto Trasformante: Adattare}

Il terzo modo \`e l'atto che trasforma. Quando una forma fallisce ---
quando la sua coerenza scende a Inerte o Fallimento --- il praticante
deve adattarsi. Deve \emph{sognare} un nuovo approccio, proporre una
variazione sul pattern fallito e sottoporla allo stesso test di coerenza.

La piattaforma \praxis\ Axiom AI serve a questa funzione: scansiona
la storia del praticante, rileva quali pattern hanno successo e quali
falliscono, e propone aggiustamenti. Questa non \`e un'imposizione
esterna ma uno specchio tenuto davanti alla storia naturale del
praticante --- un modo di \emph{vedere ci\`o che la natura sta dicendo}
sulle proprie azioni.

\section{Il Diagramma dell'Omologia}

Possiamo rappresentare l'omologia come segue:

\begin{center}
\begin{tikzpicture}[
  node distance=12mm and 24mm,
  box/.style={draw, rounded corners, align=center, minimum height=8mm,
              minimum width=28mm, font=\small},
  >={Stealth[length=2.4mm]}]
  \node[box, fill=green!10] (d-physis) {Distinzione\\$\rightarrow$\textsc{Classifica}};
  \node[box, right=of d-physis] (d-praxis) {Fare\\$\rightarrow$\textsc{Albero Obiettivi}};

  \node[box, fill=blue!10, below=of d-physis] (c-physis) {Coerenza\\$\rightarrow$\textsc{Traccia}};
  \node[box, right=of c-physis] (c-praxis) {Mantenere\\$\rightarrow$\textsc{Striscia / Check-in}};

  \node[box, fill=red!10, below=of c-physis] (v-physis) {Variazione\\$\rightarrow$\textsc{Sogno}};
  \node[box, right=of v-physis] (v-praxis) {Adattare\\$\rightarrow$\textsc{Axiom AI}};

  \node[above, font=\bfseries] at (d-physis.north) {\textsc{Physis} (Natura)};
  \node[above, font=\bfseries] at (d-praxis.north) {\textsc{Praxis} (Azione)};

  \path[<->] (d-physis) edge node[above, font=\tiny]{omologia} (d-praxis);
  \path[<->] (c-physis) edge node[above, font=\tiny]{omologia} (c-praxis);
  \path[<->] (v-physis) edge node[above, font=\tiny]{omologia} (v-praxis);
\end{tikzpicture}
\end{center}

\section{Cosa Implica l'Omologia}

Tre letture dell'omologia sono disponibili, e la tesi deve impegnarsi per una,
poich\'e non sono equivalenti e l'argomentazione che segue dipende da quale si
intenda.

La lettura \emph{debole} \`e l'\textbf{analogia}: \praxis\ semplicemente
assomiglia a \physis, come un cuore assomiglia a una pompa. La somiglianza \`e
illuminante ma inessenziale --- un sussidio didattico da scartare una volta
appresa la lezione. La lettura \emph{forte} \`e l'\textbf{identit\`a}: \praxis\
e \physis\ sono numericamente una cosa sola, sicch\'e l'azione umana \`e
semplicemente una regione della natura senza resto. La lettura \emph{intermedia}
--- quella che adottiamo --- \`e l'\textbf{omologia} nel senso biologico: non
somiglianza di superficie e non identit\`a numerica, ma \emph{lo stesso principio
generativo che produce forme a diversi livelli di organizzazione}, come l'arto
anteriore di un pipistrello e il braccio di un essere umano sono una sola
struttura diversificata dalla discendenza. L'ala e il braccio non sono analoghi
(la loro somiglianza non \`e casuale) e non sono identici (sono palesemente
organi diversi); sono omologhi --- la stessa forma sottostante, modificata dal
livello a cui si esprime.

Ci impegniamo, dunque, a questo: distinzione, coerenza e variazione non sono
metafore dell'attivit\`a umana, n\'e un insieme separato di facolt\`a umane che
per caso assomigliano a quelle della natura. Sono \emph{le tre facolt\`a proprie
della natura, ricorrenti al livello dell'azione deliberata} --- gli stessi
principi, modificati dal medium di un agente autocosciente. \praxis\ sta a
\physis\ come l'ala del pipistrello sta all'arto ancestrale: continua nel
principio, distinta nell'espressione. Questo \`e pi\`u forte dell'analogia e
pi\`u debole dell'identit\`a, e la differenza conta --- poich\'e la lettura
dell'identit\`a collasserebbe il praticante in un mero evento naturale
(dissolvendo l'agency che il resto della tesi richiede), mentre la lettura
dell'analogia renderebbe l'omologia decorativa (dissolvendo la pretesa di
governo). L'omologia preserva entrambe: l'azione \`e governata dalla natura
perch\'e \`e la stessa forma della natura, eppure resta azione --- il livello a
cui quella forma diventa responsabile verso se stessa.

\`E essenziale essere esatti su ci\`o che i sistemi funzionanti stabiliscono e
non stabiliscono. Il motore \physis\ e la piattaforma \praxis\ \emph{non}
provano l'omologia; nessun codice in esecuzione potrebbe, poich\'e l'omologia
\`e un'affermazione sulla struttura della natura, e quell'affermazione \`e
risolta dall'argomento, non dalla dimostrazione. Ci\`o che il motore mostra \`e
pi\`u ristretto e ancora decisivo: che distinzione, coerenza e variazione sono
\emph{computazionalmente trattabili} --- che possono essere costruite. La
piattaforma mostra che sono \emph{socialmente organizzabili} --- che gli esseri
umani possono strutturare le loro vite intorno a esse. La costruzione prova la
\emph{costruibilit\`a}, non la \emph{verit\`a}: rimuove l'obiezione che l'azione
governata dalla natura sia una bella impossibilit\`a, un'architettura che non
potrebbe mai reggersi. Non pu\`o, e non rimuove, l'onere dell'argomento
filosofico per l'omologia stessa. Ci\`o che la filosofia deve affermare con
l'argomento, la costruzione pu\`o solo mostrare essere \emph{non oziosa} --- un
desiderio che, almeno, \`e stato fatto funzionare.

Il resto di questa indagine esplora cosa significa questa omologia per
l'educazione del praticante (Libro III), la giustizia dell'azione
(Libro IV), l'ideale del praticante-re (Libro V) e il percorso con cui
\praxis\ degenera quando dimentica il suo fondamento in \physis\
(Libri VIII--IX).
